% Genealogía operatoria añadida en REV.2.
% Residencia única del principio de contracción de fibras y memoria.

\chapter[Arithmétique, calcul discret et mémoire généalogique]{Généalogie
opératoire de l'arithmétique et du calcul discret : composition, contraction
de fibres et relèvements avec mémoire}
\label{chap:genealogia-operaciones-rev2}
\index{opération!généalogie}
\index{mémoire!relèvement opératoire}

L'arithmétique publie des résultats ; une théorie généalogique doit conserver
aussi les conditions qui ont rendu chaque résultat possible. Cette distinction
n'ajoute pas d'opération nouvelle à l'addition, au produit, à la puissance, à la racine, au
logarithme, à la différentiation ou à l'intégration. Elle précise quelle structure
chacune contracte et quelles données un relèvement doit transporter pour distinguer
des antécédents qu'une même sortie identifie. APP montre le premier cas sur
un support commun ; TRIT oriente localement la transition ; TPK conserve la
provenance tout au long de la composition.

\section{Opération publiée et relèvement généalogique}

Soient \(X\) et \(Y\) des ensembles et soit \(f:X\to Y\) une application qui publie
une opération. Une \emph{mémoire d'opération} est un ensemble \(M_f\) muni
d'une application

\begin{equation}
 \widetilde f:X\longrightarrow Y\times M_f,
 \qquad
 \operatorname{pr}_Y\circ\widetilde f=f.
 \label{eq:levantamiento-genealogico-operacion-rev2}
\end{equation}

La première composante publie la sortie ; la seconde enregistre les données que le
lecteur décide de conserver. Le relèvement est \emph{complet sur les fibres}
si, pour tout \(y\in Y\), la restriction

\[
 \operatorname{pr}_{M_f}\circ\widetilde f\big|_{f^{-1}(y)}
 :f^{-1}(y)\longrightarrow M_f
\]

est injective. Cette condition exprime exactement que deux antécédents
ayant une même sortie visible restent distinguables par leur mémoire. Dans les
réalisations HMT, on n'exige pas qu'une unique fiche conserve toute l'histoire :
la mémoire complète peut se répartir entre le feuillet, le quotient, la retenue, la phase,
l'orientation, la frontière et le registre chronologique enrichi.

Si \(f:X\to Y\) et \(g:Y\to Z\) possèdent des relèvements fixés, leur composition
admet le relèvement canonique associé à ces relèvements,

\begin{equation}
 \widetilde{g\circ f}:X\longrightarrow Z\times M_f\times M_g,
 \qquad
 \widetilde{g\circ f}(x)
 =\bigl(g(f(x)),m_f(x),m_g(f(x))\bigr),
 \label{eq:composicion-memorias-operacion-rev2}
\end{equation}

où \(m_f\) et \(m_g\) sont les composantes de mémoire. La concaténation
n'efface pas la provenance de la première étape pour la remplacer par la seconde : elle les
coordonne. C'est le principe opératoire que le TPK réalise sur les routes.

\begin{lemma}[Composition de relèvements complets]
\label{lem:composicion-levantamientos-completos-rev2}
Si \(\widetilde f\) et \(\widetilde g\) sont complets sur les fibres de
\(f\) et de \(g\), respectivement, alors
\(\widetilde{g\circ f}\) est complet sur les fibres de \(g\circ f\).
\end{lemma}

\begin{proof}
Soient \(x_1,x_2\in X\) avec
\(g(f(x_1))=g(f(x_2))\) et avec des composantes de mémoire égales dans
\eqref{eq:composicion-memorias-operacion-rev2}. L'égalité des mémoires
\(m_g(f(x_1))=m_g(f(x_2))\), avec la complétude de
\(\widetilde g\) sur cette fibre de \(g\), implique
\(f(x_1)=f(x_2)\). L'égalité
\(m_f(x_1)=m_f(x_2)\) et la complétude de \(\widetilde f\) impliquent alors
\(x_1=x_2\).
\end{proof}

\section{Somme et produit : alternatives et composition}

Dans une algèbre de chemins, la somme réunit des alternatives aux extrémités
compatibles,

\[
 a_1\gamma_1+a_2\gamma_2,
\]

tandis que le produit concatène des étapes lorsque l'extrémité de la première
coïncide avec l'origine de la seconde,

\[
 \gamma_2\gamma_1,
 \qquad t(\gamma_1)=s(\gamma_2).
\]

La somme conserve la coexistence linéaire des histoires ; le produit construit
une histoire composée et, lorsque seule sa valeur est publiée, peut oublier la
factorisation qui l'a produite. Sur la carte APP, les évaluations
\(\Sigma\) et \(\Pi\) agissent sur le même état \((i,j)\), mais organisent des fibres
distinctes. Dans le bloc central \(B_c\) d'APP, la séparation spectrale qui
est développée dans le chapitre suivant mesure exactement cette incompatibilité
locale ; ce n'est pas un écart entre deux tables indépendantes et elle n'est pas présentée
comme une distance globale entre tous les lecteurs.

Dans le feuillet additif, la translation \(T_a(x)=x+a\) est réversible sur le groupe
résiduel. Dans le feuillet multiplicatif, \(M_a(x)=ax\) est inversible dans
\(\mathbb Z/9\mathbb Z\) exactement lorsque \(\gcd(a,9)=1\). Les classes non
unitaires \(\overline0,\overline3,\overline6\), publiées dans la carte
positive comme \(9,3,6\), replient les fibres et rendent visible la nécessité de
conserver le quotient et la provenance. La somme et le produit constituent ainsi les deux
lecteurs inauguraux d'une même géométrie des chemins.

\begin{lemma}[Conjugaison affine dans le secteur unitaire]
\label{lem:conjugacion-afin-app-rev3}
Soit \(R=\mathbb Z/9\mathbb Z\). Pour \(a\in R^\times\) et \(b\in R\), les
applications \(T_b(x)=x+b\) et \(M_a(x)=ax\) satisfont
\begin{equation}
 M_aT_bM_a^{-1}=T_{ab}.
 \label{eq:conjugacion-afin-app-rev3}
\end{equation}
Si \(a\in(3)\), \(M_a\) cesse d'être inversible et la conjugaison n'est pas
définie : son image reste contenue dans le radical et plusieurs branches du feuillet
additif sont identifiées.
\end{lemma}

\begin{proof}
Pour \(a\in R^\times\),
\(
 M_aT_bM_a^{-1}(x)=a(a^{-1}x+b)=x+ab=T_{ab}(x)
\).
Si \(a\in(3)\), le noyau de \(M_a\) est non trivial ; il n'existe donc pas de
\(M_a^{-1}\) et l'expression de conjugaison est dépourvue de type. Le cas
\(a=3\) réalise explicitement le repliement au moyen de
\(\ker M_3=\operatorname{im}M_3=(3)\).
\end{proof}

\section{Puissance, logarithme et racine}

Soit \(x\) un élément d'un monoïde, ou d'une algèbre associative, et soit
\(n\geq1\). La puissance enregistre l'itération d'une composition :

\[
 x^n=\underbrace{x\cdots x}_{n\ \mathrm{factores}}.
\]

L'exposant est la profondeur de cette itération et l'orbite de \(x\) détermine
quelles sorties sont identifiées. Dans le domaine réel positif, si \(b>0\),
\(b\neq1\) et \(x>0\), on a
\[
 \log_b(x^n)=n\log_b x.
\]
Dans le domaine complexe, la même identité requiert une branche compatible et
que l'itération ne franchisse pas sa coupure. Le logarithme donne une coordonnée à la profondeur
\(n\) seulement après avoir également fixé le générateur \(x\) et exigé
\(\log_b x\neq0\). Base, générateur, domaine, branche et coupure font partie du
lecteur.

La racine inverse une puissance comme problème de fibres,

\[
 y^n=x.
\]

La fibre peut être vide, unitaire ou multiple. Sélectionner une racine requiert des
données de branche, d'orientation, de phase ou de frontière. La construction ultérieure d'un
opérateur interne \(J_E\in\operatorname{End}(V)\), avec \(J_E^2=-I_V\),
illustre ce principe : le
symbole \(\sqrt{-1}\) publie une coordonnée, tandis que l'opérateur, son domaine
et son action conservent la structure qui la soutient.

\section{Calcul discret sur des cylindres compatibles : opérateurs de différentiation et d'intégration}

Soit \(P_n\) un chemin orienté connexe et soit \(A\) un groupe abélien. Dans le
complexe discret de mesure, le cobord combinatoire

\[
 d:C^0(P_n;A)\longrightarrow C^1(P_n;A),
 \qquad (df)(e)=f(t(e))-f(s(e)),
\]

a pour noyau exactement les \(0\)-cochaînes constantes. Reconstruire
\(f\) à partir de \(df\) exige une condition au bord, par exemple la valeur en un
sommet de référence.

Le cobord dépend de l'incidence et de l'orientation, mais non d'une
longueur. Si \(A\) est en outre un espace vectoriel réel et
\(\ell:E(P_n)\to\mathbb R_{>0}\) une longueur d'arête, la dérivée
métrique est définie par
\[
 D_\ell f(e)=\frac{(df)(e)}{\ell(e)}.
\]
Pour une route \(\gamma=(e_1,\ldots,e_r)\) compatible avec les orientations,
\begin{equation}
 \sum_{k=1}^r D_\ell f(e_k)\,\ell(e_k)
 =f(t(\gamma))-f(s(\gamma)).
 \label{eq:integracion-derivada-metrica-rev2}
\end{equation}
L'intégration accumule ainsi une cochaîne le long d'une route et récupère la
différence des extrémités lorsque la cochaîne est exacte. Le cobord, la dérivée métrique
et l'intégrale dépendent de données distinctes : incidence et orientation ; longueur et
échelle ; chemin et section de bord.

Cette dépendance explique pourquoi le passage du comptage des marques à la mesure des intervalles
est structurel. La dérivation vit sur les relations interstitielles ;
l'intégration compose ces relations ; le cycle ajoute le porteur topologique de
l'holonomie ; le transport TPK détermine quelle mémoire survit au tour.

\section{Hiérarchie de fibres et fonction du TPK}

La hiérarchie peut être condensée sans identifier ses niveaux :

\begin{center}
\small
\begin{tabularx}{\textwidth}{@{}
  >{\raggedright\arraybackslash}p{0.13\textwidth}
  >{\raggedright\arraybackslash}p{0.22\textwidth}
  >{\raggedright\arraybackslash}X
  >{\raggedright\arraybackslash}p{0.25\textwidth}@{}}
\toprule
\textbf{Opération} & \textbf{Action première} &
\textbf{Fibre contractée par la sortie} & \textbf{Mémoire de reconstruction}\\
\midrule
Somme & agrégation ou translation & alternatives ayant une même évaluation &
origine, coefficients et support\\
Produit & composition ou échelle & factorisations et classes non unitaires &
ordre des facteurs, quotient et feuillet\\
Puissance & itération & orbites de composition &
exposant, ordre et point initial\\
Racine & inversion multivoque & antécédents d'une puissance &
branche, phase, orientation et frontière\\
Logarithme & coordonnée de profondeur & périodes et déterminations &
base, générateur, domaine, branche et coupure\\
Cobord et dérivée métrique & différence locale et normalisation & modo constante &
incidence, orientation, longueur, échelle et référence\\
Intégrale & accumulation de chemin & routes ayant la même valeur publiée &
chemin, mesure, bord et holonomie\\
\bottomrule
\end{tabularx}
\end{center}

TPK conserve ces coordonnées à l'intérieur d'états transportables et non comme une
liste ajoutée au résultat. Sa fiche \((D,C,f,I,L,M)\) déclare pour chaque
opérateur le domaine, le codomaine, la règle, les invariants, l'information
publiée ou perdue et la mémoire conservée. La transition

\[
 \text{opération visible}
 \longrightarrow
 \text{relèvement avec mémoire}
 \longrightarrow
 \text{section compatible}
\]

fournit le lien entre l'arithmétique inaugurale, la reconstruction du
continu et la définition ultérieure du nombre généalogique. La frontière
nombre--particule conservera la même architecture en ajoutant une réalisation
dimensionnelle à la section survivante.

\paragraph{Domaine inaugural de la cellule nonadique et exclusion causale de l’APP, du TRIT, du TPK et des constantes analytiques}
La priorité de la somme et du produit dans APP, la conservation de la phase, du feuillet,
de la retenue et de la frontière par TPK, et la transition de la section à la route constituent les
prémisses du relèvement
\eqref{eq:levantamiento-genealogico-operacion-rev2}, de la notion de complétude
sur les fibres et de la table des mémoires. On n'identifie pas une opération à sa réalisation physique :
chaque étape ultérieure conserve son domaine, sa carte et sa source.
